%% ============================================================
%%  RAPPORT DE MÉMOIRE — LICENCE PRO (BAC+3)
%%  Plateforme Woyofal — Analyse de Données et Simulation de Recharges Électriques
%%  Auteur : Diarra NDOYE — Sonatel Academy (ECSA) — Avril 2026
%% ============================================================
\documentclass[12pt,a4paper]{report}

%% ── Encodage & langue ───────────────────────────────────────
\usepackage[utf8]{inputenc}
\usepackage[T1]{fontenc}
\usepackage[french]{babel}

%% ── Mise en page ────────────────────────────────────────────
\usepackage[top=2.5cm,bottom=2.5cm,left=3cm,right=2.5cm]{geometry}
\usepackage{setspace}
\onehalfspacing
\setlength{\parindent}{0pt}
\setlength{\parskip}{6pt}

%% ── Polices ─────────────────────────────────────────────────
\usepackage{lmodern}
\usepackage{microtype}

%% ── Couleurs & graphiques ───────────────────────────────────
\usepackage[dvipsnames]{xcolor}
\definecolor{rouge}{HTML}{DC2626}
\definecolor{marine}{HTML}{1A1D2E}
\definecolor{grislight}{HTML}{F3F4F6}
\definecolor{grisborder}{HTML}{E5E7EB}
\usepackage{graphicx}
\usepackage{tikz}
\usetikzlibrary{shapes.geometric, arrows.meta, positioning, fit, backgrounds}

%% ── Tableaux ────────────────────────────────────────────────
\usepackage{booktabs}
\usepackage{array}
\usepackage{longtable}
\usepackage{tabularx}
\usepackage{multirow}
\usepackage{colortbl}

%% ── Maths ───────────────────────────────────────────────────
\usepackage{amsmath}

%% ── Code source ─────────────────────────────────────────────
\usepackage{listings}
\lstset{
  backgroundcolor=\color{grislight},
  basicstyle=\ttfamily\small,
  keywordstyle=\color{rouge}\bfseries,
  commentstyle=\color{gray},
  stringstyle=\color{marine},
  breaklines=true,
  frame=single,
  rulecolor=\color{grisborder},
  numbers=left,
  numberstyle=\tiny\color{gray},
  numbersep=8pt,
  captionpos=b,
}
\lstdefinelanguage{JavaScript}{
  keywords={const,let,var,function,return,import,export,default,
            from,if,else,for,while,async,await,try,catch,new,class},
  morestring=[b]",
  morestring=[b]',
  morestring=[b]`,
  morecomment=[l]{//},
  morecomment=[s]{/*}{*/},
  sensitive=true
}

%% ── En-têtes & pieds de page ────────────────────────────────
\usepackage{fancyhdr}
\pagestyle{fancy}
\fancyhf{}
\fancyhead[L]{\small\textcolor{marine}{\textit{Plateforme Woyofal}}}
\fancyhead[R]{\small\textcolor{rouge}{Diarra NDOYE}}
\fancyfoot[C]{\thepage}
\renewcommand{\headrulewidth}{0.4pt}

%% ── Liens hypertexte ────────────────────────────────────────
\usepackage[colorlinks=true,linkcolor=rouge,urlcolor=rouge,citecolor=marine]{hyperref}

%% ── Listes personnalisées ───────────────────────────────────
\usepackage{enumitem}
\setlist{noitemsep, topsep=4pt}

%% ── Boîtes encadrées ────────────────────────────────────────
\usepackage{tcolorbox}
\tcbuselibrary{skins,breakable}
\newtcolorbox{infobox}[1][]{
  colback=grislight, colframe=rouge, coltitle=white,
  fonttitle=\bfseries\small, title=#1,
  left=6pt, right=6pt, top=4pt, bottom=4pt,
  arc=4pt, boxrule=0.8pt
}
\newtcolorbox{codebox}[1][]{
  colback=marine!5, colframe=marine!30, coltitle=marine,
  fonttitle=\ttfamily\small, title=#1,
  left=6pt, right=6pt, top=4pt, bottom=4pt,
  arc=4pt, boxrule=0.8pt
}

%% ── Titre de chapitre personnalisé ──────────────────────────
\usepackage{titlesec}
\titleformat{\chapter}[block]
  {\normalfont\huge\bfseries\color{marine}}
  {\textcolor{rouge}{\thechapter.}}{12pt}{}
  [\vspace{2pt}\rule{\textwidth}{1.5pt}]
\titlespacing*{\chapter}{0pt}{20pt}{20pt}
\titleformat{\section}{\large\bfseries\color{marine}}{\textcolor{rouge}{\thesection}}{8pt}{}
\titleformat{\subsection}{\normalsize\bfseries\color{marine!80}}{\thesubsection}{6pt}{}

%% ── Métadonnées PDF ─────────────────────────────────────────
\hypersetup{
  pdftitle={Rapport de Mémoire — Plateforme Woyofal},
  pdfauthor={Diarra NDOYE},
  pdfsubject={Analyse de Données et Simulation de Recharges Électriques — Sénégal},
  pdfkeywords={Woyofal, Senelec, ETL, Data Warehouse, Machine Learning, React, simulation}
}

%% ============================================================
\begin{document}

%% ── PAGE DE GARDE ───────────────────────────────────────────
\begin{titlepage}
  \centering
  \vspace*{1.5cm}

  \begin{tikzpicture}
    \fill[marine] (0,0) rectangle (\textwidth,1.2cm);
    \node[white,font=\Large\bfseries] at (0.5\textwidth,0.6cm)
      {École du Code Sonatel Academy (ECSA) — Orange Digital Center};
  \end{tikzpicture}

  \vspace{1.2cm}
  {\large Licence Professionnelle — Développement Data}\\[0.3cm]
  {\small Référentiel Dev Data — Promotion 7 (Cohorte 6)}

  \vspace{1.5cm}
  \rule{\textwidth}{0.8pt}\\[0.4cm]
  {\Huge\bfseries\color{marine} Plateforme Woyofal}\\[0.3cm]
  {\Large\color{rouge} Conception et réalisation d'une Plateforme\\
  pour l'Analyse de Données et Simulation\\de Recharges Électriques}\\
  {\large au Sénégal}\\[0.4cm]
  \rule{\textwidth}{0.8pt}

  \vspace{1.5cm}
  \begin{tabular}{rl}
    \textbf{Auteur :}        & Diarra NDOYE \\[4pt]
    \textbf{Encadrant :}     & M. Mamadou MBAYE \\[4pt]
    \textbf{Site web :}      & \href{https://frontend-react-wine.vercel.app}{frontend-react-wine.vercel.app} \\[4pt]
    \textbf{Dépôt GitHub :}  & \href{https://github.com/ndoye09/woyofal2026}{github.com/ndoye09/woyofal2026} \\[4pt]
    \textbf{Établissement :} & ECSA / Orange Digital Center \\[4pt]
    \textbf{Année :}         & 2025--2026 \\
  \end{tabular}

  \vfill
  {\small Dakar, Sénégal — Avril 2026}
\end{titlepage}

%% ── RÉSUMÉ ──────────────────────────────────────────────────
\chapter*{Résumé}
\addcontentsline{toc}{chapter}{Résumé}

Ce mémoire présente la conception et le déploiement de \textbf{Woyofal}, une plateforme complète d'ingénierie des données et de simulation de recharges électriques prépayées au Sénégal, basée sur la grille tarifaire officielle Senelec 2026 (Décision CRSE~2025-140).

Le projet couvre l'intégralité de la chaîne de valeur de la donnée : génération d'un jeu de données synthétiques de \textbf{330~000 lignes}, pipeline ETL orchestré par \textbf{Apache Airflow}, entrepôt de données \textbf{PostgreSQL} en schéma en étoile (4 dimensions, 2 tables de faits), transformations \textbf{dbt}, analyse exploratoire (EDA) complète, modèle de \textbf{Machine Learning} (Random Forest, accuracy 99,75\%), et une interface web React~18 déployée en production sur Vercel.

La couche frontend intègre un \textbf{simulateur bidirectionnel} (FCFA $\rightarrow$ kWh et kWh $\rightarrow$ FCFA), un mini-calculateur interactif, un assistant IA spécialisé (OpenRouter, cascade Llama~3.3~70B / Gemma~3 / Qwen~3 / Mistral~7B), ainsi qu'un historique des recharges.

\textbf{Mots-clés :} Woyofal, Senelec, prépayé, tarifs 2026, ETL, Data Warehouse, PostgreSQL, dbt, Airflow, Machine Learning, React, Vercel, simulation tarifaire.

\vspace{1cm}
\textbf{Abstract} \\[4pt]
This thesis presents the design and deployment of \textit{Woyofal}, a full data engineering and simulation platform for prepaid electricity recharges in Senegal, based on the official Senelec 2026 tariff grid (CRSE Decision 2025-140). The project covers the entire data value chain: synthetic data generation (330,000 rows), an Airflow ETL pipeline, a PostgreSQL star-schema data warehouse, dbt transformations, exploratory data analysis, a Machine Learning model (Random Forest, 99.75\% accuracy), and a React 18 frontend deployed on Vercel.

\newpage

%% ── TABLE DES MATIÈRES ──────────────────────────────────────
\tableofcontents
\newpage

%% ── LISTE DES FIGURES ───────────────────────────────────────
\listoffigures
\addcontentsline{toc}{chapter}{Liste des figures}
\newpage

%% ── LISTE DES TABLEAUX ──────────────────────────────────────
\listoftables
\addcontentsline{toc}{chapter}{Liste des tableaux}
\newpage

%% ── LISTE DES ABRÉVIATIONS ──────────────────────────────────
\chapter*{Liste des abréviations}
\addcontentsline{toc}{chapter}{Liste des abréviations}

\begin{tabular}{ll}
\toprule
\textbf{Abréviation} & \textbf{Définition} \\
\midrule
DWH   & Data Warehouse (entrepôt de données) \\
ETL   & Extract, Transform, Load \\
EDA   & Exploratory Data Analysis \\
FCFA  & Franc CFA (monnaie d'Afrique de l'Ouest) \\
API   & Application Programming Interface \\
CI/CD & Continuous Integration / Continuous Deployment \\
ECSA  & École du Code Sonatel Academy \\
ODC   & Orange Digital Center \\
SENELEC & Société Nationale d'Électricité du Sénégal \\
CRSE  & Commission de Régulation du Secteur de l'Énergie \\
DPP   & Domestique Petite Puissance (type compteur) \\
PPP   & Professionnel Petite Puissance (type compteur) \\
SPA   & Single Page Application \\
JWT   & JSON Web Token \\
ML    & Machine Learning \\
\bottomrule
\end{tabular}
\newpage

%% ============================================================
\chapter{Introduction}
%% ============================================================

\section{Contexte général}

Dans un Sénégal en pleine transformation numérique, l'accès à l'électricité constitue un pilier fondamental du développement économique et social. Le système de recharges prépayées, communément appelé \textit{Woyofal}, est au cœur du quotidien de millions de ménages sénégalais.

Le Sénégal compte plus de 17 millions d'habitants répartis dans 14 régions et 23 zones d'approvisionnement électrique gérées par la \textbf{Senelec} (Société Nationale d'Électricité du Sénégal). Le compteur prépayé \textit{Woyofal} est commercialisé par Senelec depuis 2012 : l'utilisateur achète des tokens chez un revendeur ou via mobile money, les saisit sur le compteur, et le crédit est converti en kWh selon la grille tarifaire en vigueur.

Le système tarifaire sénégalais, régi par la décision \textbf{CRSE 2025-140}, applique un barème progressif à plusieurs tranches selon la consommation mensuelle cumulée. Cette progressivité — bien que socialement justifiée — est d'une complexité opaque pour l'utilisateur final : combien de kWh obtiendra-t-il pour une recharge de 5~000~FCFA ? Va-t-il franchir le seuil de la tranche~1 ?

C'est dans ce contexte que s'inscrit ce mémoire, fruit d'une formation rigoureuse au sein de la \textbf{Sonatel Academy}, dans le cadre du référentiel Développement Data (Promo~7 / Cohorte~6).

\section{Problématique}

\begin{infobox}[Problème identifié]
Les consommateurs Woyofal n'ont \textbf{aucun outil simple et gratuit} pour simuler leurs recharges selon les tarifs Senelec 2026 officiels et analyser leurs comportements de consommation. Les outils existants utilisent des tarifs obsolètes, ignorent le cumul mensuel (code~814) et ne proposent aucune infrastructure de données pour l'analyse à l'échelle nationale.
\end{infobox}

De façon plus précise, les problèmes soulevés sont :
\begin{itemize}
  \item Absence d'outil de simulation intégrant la grille tarifaire 2026 actualisée (Décision 2025-140) ;
  \item Incompréhension du mécanisme des tranches et de la redevance mensuelle (429~FCFA) ;
  \item Impossibilité de planifier son budget électrique à partir d'une cible en kWh (\textit{calcul inverse}) ;
  \item Données de consommation non exploitées faute d'entrepôt et de pipeline ETL ;
  \item Aucune visualisation des disparités de consommation entre les 23 zones géographiques.
\end{itemize}

\section{Objectifs}

\subsection*{Objectif général}
Concevoir et implémenter une plateforme d'analyse de données complète, intégrant un pipeline ETL robuste, un Data Warehouse structuré, un simulateur de recharges conforme à la grille SENELEC 2026, et des interfaces de visualisation interactives accessibles à tous.

\subsection*{Objectifs spécifiques}
\begin{enumerate}
  \item \textbf{O1} — Générer et structurer un jeu de données réaliste de 330~000+ lignes représentatif des comportements au Sénégal ;
  \item \textbf{O2} — Mettre en place un pipeline ETL robuste (Airflow) vers un Data Warehouse PostgreSQL en star schema ;
  \item \textbf{O3} — Réaliser une analyse exploratoire (EDA) complète des comportements de consommation et de recharge ;
  \item \textbf{O4} — Entraîner un modèle de Machine Learning capable de prédire le dépassement de tranche~1 ;
  \item \textbf{O5} — Implémenter le moteur de calcul tarifaire DPP/PPP 2026 (4~tranches, redevance, taxe) avec mode inverse (kWh $\rightarrow$ FCFA) ;
  \item \textbf{O6} — Déployer une interface React accessible, responsive, avec CI/CD GitHub $\rightarrow$ Vercel ;
  \item \textbf{O7} — Intégrer un assistant IA conversationnel pour l'aide tarifaire.
\end{enumerate}

\section{Structure du rapport}

Le présent rapport s'organise comme suit : le chapitre~\ref{ch:contexte} présente la grille tarifaire, le marché Woyofal et la critique de l'existant. Le chapitre~\ref{ch:archi} décrit l'architecture technique globale de la plateforme. Le chapitre~\ref{ch:data} présente le pipeline de données, l'entrepôt PostgreSQL et les transformations dbt. Le chapitre~\ref{ch:eda_ml} expose l'analyse exploratoire des données et le modèle de Machine Learning. Le chapitre~\ref{ch:impl} détaille l'implémentation du simulateur et de l'interface frontend. Le chapitre~\ref{ch:resultats} présente les résultats et tests. Le chapitre~\ref{ch:conclusion} conclut et ouvre des perspectives.

%% ============================================================
\chapter{Contexte et Problématique}
\label{ch:contexte}
%% ============================================================

\section{Le système prépayé Woyofal}

\subsection*{Fonctionnement du cumul mensuel}
À chaque début de mois, le compteur réinitialise le \textbf{cumul mensuel} à zéro. Ce cumul (consultable via le code~814 sur le compteur) détermine la tranche tarifaire applicable à chaque recharge. La première recharge du mois inclut une \textbf{redevance fixe de 429~FCFA} pour les compteurs DPP (location et entretien compteur).

\section{Grille tarifaire Senelec 2026}

La décision CRSE 2025-140 définit deux catégories tarifaires pour les compteurs prépayés.

\subsection{Tarif DPP — Domestique Petite Puissance}

\begin{table}[h!]
\centering
\caption{Tarifs DPP 2026 — Domestique Petite Puissance}
\label{tab:dpp}
\begin{tabularx}{0.85\textwidth}{>{\bfseries}l c X c}
\toprule
\rowcolor{marine!10}
Tranche & Seuil (kWh/mois) & Description & Prix (FCFA/kWh) \\
\midrule
T1 & 0 -- 150    & Tarif social (baisse 10\% vs 2024) & \textbf{82,00} \\
T2 & 151 -- 250  & Tranche intermédiaire & \textbf{136,49} \\
T3 & 251 -- 400  & Tranche haute         & \textbf{136,49} \\
T4 & $>$ 400     & Très haute consommation & \textbf{162,00} \\
\midrule
\multicolumn{3}{l}{Redevance mensuelle (1\textsuperscript{ère} recharge)} & \textbf{429 FCFA} \\
\multicolumn{3}{l}{Taxe communale} & \textbf{2,5\%} \\
\bottomrule
\end{tabularx}
\end{table}

\subsection{Tarif PPP — Professionnel Petite Puissance}

\begin{table}[h!]
\centering
\caption{Tarifs PPP 2026 — Professionnel Petite Puissance}
\label{tab:ppp}
\begin{tabularx}{0.82\textwidth}{>{\bfseries}l c X c}
\toprule
\rowcolor{marine!10}
Tranche & Seuil (kWh/mois) & Description & Prix (FCFA/kWh) \\
\midrule
T1 & 0 -- 50    & Tarif de base            & \textbf{147,43} \\
T2 & 51 -- 500  & Tranche standard         & \textbf{189,84} \\
T3 & $>$ 500    & Grande consommation      & \textbf{189,84} \\
\bottomrule
\end{tabularx}
\end{table}

\section{Critique de l'existant}

\subsection{Solution de référence : woyofal-check.site}

Le site \textit{woyofal-check.site} constitue la solution en ligne la plus avancée disponible pour les consommateurs Woyofal. Développé en 2025, il propose un calculateur, un suivi et des prédictions basiques. Cependant, une analyse approfondie révèle plusieurs limitations qui justifient le développement d'une plateforme plus robuste.

\textbf{Limites fonctionnelles identifiées :}
\begin{itemize}
  \item Accès conditionné à la création d'un compte pour toutes les fonctionnalités utiles ;
  \item Grille tarifaire antérieure à la Décision 2025-140 (prix obsolètes) ;
  \item Absence de prise en compte du cumul mensuel (code 814) ;
  \item Aucun enregistrement des simulations ;
  \item Pas de chatbot — FAQ statique uniquement ;
  \item Bugs de connexion constatés.
\end{itemize}

\textbf{Limites architecturales :}
\begin{itemize}
  \item Aucun Data Warehouse ni pipeline ETL ;
  \item Pas de dashboard analytique ;
  \item Pas d'API REST documentée ou réutilisable.
\end{itemize}

\subsection{Tableau comparatif}

\begin{table}[h!]
\centering
\caption{Comparaison Plateforme Woyofal vs solution existante}
\label{tab:comparatif}
\begin{tabularx}{\textwidth}{X c c}
\toprule
\rowcolor{marine!10}
Critère & \textbf{Woyofal (ce projet)} & woyofal-check.site \\
\midrule
Calculateur FCFA $\rightarrow$ kWh & \textcolor{green!60!black}{\textbf{✓}} & \textcolor{green!60!black}{\textbf{✓}} \\
Calcul inverse kWh $\rightarrow$ FCFA (sans connexion) & \textcolor{green!60!black}{\textbf{✓}} & \textcolor{orange}{\textbf{partiel}} \\
Grille tarifaire 2026 officiell (4 tranches DPP) & \textcolor{green!60!black}{\textbf{✓}} & \textcolor{red}{\textbf{✗}} \\
Prise en compte cumul mensuel (code 814) & \textcolor{green!60!black}{\textbf{✓}} & \textcolor{red}{\textbf{✗}} \\
Historique des simulations & \textcolor{green!60!black}{\textbf{✓}} & \textcolor{red}{\textbf{✗}} \\
Chatbot IA (OpenRouter multi-modèles) & \textcolor{green!60!black}{\textbf{✓}} & \textcolor{red}{\textbf{✗}} \\
Data Warehouse / pipeline ETL & \textcolor{green!60!black}{\textbf{✓}} & \textcolor{red}{\textbf{✗}} \\
Dashboard analytique (23 zones) & \textcolor{green!60!black}{\textbf{✓}} & \textcolor{red}{\textbf{✗}} \\
Déploiement cloud CI/CD & \textcolor{green!60!black}{\textbf{✓}} Vercel & \textcolor{orange}{\textbf{PWA basique}} \\
\bottomrule
\end{tabularx}
\end{table}

\section{Analyse du besoin}

Une enquête informelle auprès de ménages dakarois (Janvier 2026) met en évidence :
\begin{itemize}
  \item 87\% des utilisateurs ignorent leur cumul mensuel actuel avant de recharger ;
  \item La quasi-totalité ne connaît pas le montant exact pour rester en tranche~1 ;
  \item Aucun outil gratuit ne gère correctement les 4 tranches DPP 2026.
\end{itemize}

%% ============================================================
\chapter{Architecture Technique}
\label{ch:archi}
%% ============================================================

\section{Vue d'ensemble}

La plateforme Woyofal est organisée en deux grandes couches complémentaires : une \textbf{couche data} (ingénierie des données, entrepôt, Machine Learning) et une \textbf{couche présentation} (frontend React, simulateur, assistant IA). L'ensemble s'appuie sur une architecture \textbf{JAMstack}\footnote{JavaScript, API, Markup — architecture découplant frontend et backend.} pour maximiser la disponibilité mobile et minimiser les coûts d'infrastructure.

\begin{figure}[h!]
\centering
\begin{tikzpicture}[
  box/.style={draw=marine, fill=marine!8, rounded corners=4pt,
              minimum width=3.0cm, minimum height=0.85cm,
              font=\small\bfseries, text=marine, align=center},
  boxred/.style={draw=rouge, fill=rouge!8, rounded corners=4pt,
              minimum width=3.0cm, minimum height=0.85cm,
              font=\small\bfseries, text=rouge, align=center},
  arrow/.style={-{Latex[length=6pt]}, thick, color=marine!60},
  node distance=1.1cm and 1.5cm
]
  \node[box] (gen) {Génération CSV\\(330k lignes)};
  \node[box, right=of gen] (etl) {ETL Airflow\\(Pipeline DAG)};
  \node[box, right=of etl] (pg) {PostgreSQL\\(Star Schema)};
  \node[box, right=of pg] (dbt) {dbt\\(Transformations)};
  \node[box, right=of dbt] (ml) {ML RandomForest\\(99,75\% acc.)};

  \node[box, below=1.8cm of pg] (react) {Frontend React 18\\+ Vite + Tailwind};
  \node[box, left=of react] (api) {API Mock\\Node.js :5000};
  \node[boxred, right=of react] (vercel) {Vercel CDN\\(Production)};
  \node[box, below=of react] (user) {Utilisateur\\(Mobile / Web)};

  \draw[arrow] (gen) -- (etl);
  \draw[arrow] (etl) -- (pg);
  \draw[arrow] (pg) -- (dbt);
  \draw[arrow] (dbt) -- (ml);
  \draw[arrow] (pg) -- (react);
  \draw[arrow] (api) -- node[above,font=\tiny]{dev} (react);
  \draw[arrow] (react) -- (vercel);
  \draw[arrow] (react) -- (user);
\end{tikzpicture}
\caption{Architecture complète de la plateforme Woyofal — couche data et couche présentation}
\label{fig:archi}
\end{figure}

\section{Architecture globale — 7 couches}

\begin{table}[h!]
\centering
\caption{Architecture en 7 couches de la plateforme}
\label{tab:couches}
\begin{tabularx}{\textwidth}{c l X}
\toprule
\rowcolor{marine!10}
Couche & Nom & Description \\
\midrule
1 & Sources         & CSV : zones, users, consumption\_daily, recharges \\
2 & Ingestion       & Validation (nulls, plages, intégrité), retry, logging \\
3 & Transformation  & Calcul tranches T1/T2/T3/T4, économies, agrégations \\
4 & Data Warehouse  & Star Schema Kimball : 4 dimensions + 2 tables de faits \\
5 & Calculs avancés & Simulateur recharges, prédictions ML, rapports \\
6 & Analytics       & Streamlit, Jupyter EDA, Plotly, Folium \\
7 & CI/CD           & GitHub Actions, pytest, Docker Compose, Vercel \\
\bottomrule
\end{tabularx}
\end{table}

\section{Stack technologique}

\begin{table}[h!]
\centering
\caption{Stack technologique complet}
\label{tab:stack}
\begin{tabularx}{\textwidth}{l l X}
\toprule
\rowcolor{marine!10}
Couche & Technologie & Justification \\
\midrule
\multicolumn{3}{l}{\textit{Couche données}} \\
Génération  & Python (Faker, NumPy)       & Données synthétiques réalistes, reproductibles \\
Orchestration & Apache Airflow            & DAG quotidien, retry automatique, monitoring \\
Base de données & PostgreSQL 14           & Robustesse, SQL avancé, intégration dbt \\
Conteneur   & Docker Compose              & Reproductibilité, isolation des services \\
Transformations & dbt                     & SQL versionné, tests qualité, lineage \\
ML          & scikit-learn (RandomForest) & Robustesse sur données tabulaires \\
EDA         & Pandas, Matplotlib, Seaborn & Analyse exploratoire et visualisations \\
\midrule
\multicolumn{3}{l}{\textit{Couche présentation}} \\
UI          & React 18                    & Réactivité, composants réutilisables \\
Build       & Vite 5                      & HMR ultra-rapide, build optimisé (ESBuild) \\
Style       & Tailwind CSS 3.4            & Utility-first, pas de CSS mort en production \\
Graphiques  & Recharts                    & Natif React, responsive \\
Routing     & React Router v6             & SPA avec routes protégées \\
Chatbot IA  & OpenRouter                  & Cascade Llama 3.3 70B / Gemma 3 / Qwen 3 / Mistral 7B \\
Deploy      & Vercel                      & CI/CD gratuit, CDN mondial, HTTPS auto \\
\bottomrule
\end{tabularx}
\end{table}

%% ============================================================
\chapter{Plateforme de Données}
\label{ch:data}
%% ============================================================

\section{Génération des données synthétiques}

En l'absence de données réelles accessibles, un jeu de données synthétiques a été conçu pour refléter fidèlement les comportements de consommation électrique au Sénégal. La génération est pilotée par le script \texttt{scripts/generate\_all\_data.py} qui orchestre quatre modules séquentiels.

\begin{table}[h!]
\centering
\caption{Jeu de données synthétiques généré}
\label{tab:datasets}
\begin{tabularx}{\textwidth}{l r X}
\toprule
\rowcolor{marine!10}
Fichier CSV & Volume & Contenu \\
\midrule
\texttt{zones\_senegal.csv}     & 23 lignes     & 23 zones géographiques, 14 régions, type urbain/rural \\
\texttt{users.csv}              & 10~000 lignes & Abonnés avec type compteur DPP/PPP, zone, conso moyenne \\
\texttt{consumption\_daily.csv} & 300~000 lignes & Conso quotidienne, tranche, coût FCFA, cumul mensuel \\
\texttt{recharges.csv}          & 20~000 lignes & Transactions avec montant, kWh obtenus, tranche \\
\midrule
\textbf{Total} & \textbf{330~023 lignes} & \multicolumn{1}{c}{---} \\
\bottomrule
\end{tabularx}
\end{table}

Chaque enregistrement de consommation intègre une colonne \texttt{economie\_baisse\_10pct} calculée selon la Décision 2025-140 : la baisse tarifaire de 10\% en T1 (82~FCFA/kWh au lieu de 91~FCFA antérieurement) s'applique exclusivement à la tranche sociale.

\section{Pipeline ETL avec Apache Airflow}

\subsection{Architecture du DAG}

Le pipeline ETL est orchestré par un DAG Airflow \texttt{woyofal\_etl\_pipeline} déclenché quotidiennement à 2h du matin. Il est structuré en quatre groupes de tâches :

\begin{enumerate}
  \item \textbf{check\_prerequisites} — Vérification de la disponibilité de la base et des dossiers requis ;
  \item \textbf{extract} — Chargement des CSV depuis \texttt{data/raw/} avec validation de schéma ;
  \item \textbf{transform} — Nettoyage (doublons, nulls, encodages), enrichissement tarifaire, normalisation ;
  \item \textbf{load} — Ingestion dans PostgreSQL via \texttt{COPY} server-side, puis alimentation des dimensions et faits.
\end{enumerate}

\begin{codebox}[Configuration du DAG Airflow]
\begin{lstlisting}[language=Python, numbers=none]
dag = DAG(
    'woyofal_etl_pipeline',
    default_args={
        'retries': 2,
        'retry_delay': timedelta(minutes=5),
        'email_on_failure': True,
    },
    schedule_interval='0 2 * * *',   # 2h du matin
    catchup=False,
    tags=['woyofal', 'etl', 'production']
)
\end{lstlisting}
\end{codebox}

\section{Entrepôt de données PostgreSQL}

\subsection{Schéma en étoile}

L'entrepôt de données adopte un \textbf{schéma en étoile} (star schema, méthodologie Kimball) composé de 4 tables de dimensions et 2 tables de faits.

\begin{figure}[h!]
\centering
\begin{tikzpicture}[
  dim/.style={draw=marine, fill=marine!8, rounded corners=3pt,
              minimum width=3.0cm, minimum height=1.3cm,
              font=\small\bfseries, text=marine, align=center},
  fact/.style={draw=rouge, fill=rouge!8, rounded corners=3pt,
               minimum width=3.2cm, minimum height=1.3cm,
               font=\small\bfseries, text=rouge, align=center},
  arr/.style={-{Latex[length=5pt]}, rouge, thick},
  node distance=2.0cm
]
  \node[fact] (fr) {fact\_recharges\\\footnotesize(20k lignes)};
  \node[fact, right=2.5cm of fr] (fc) {fact\_consumption\\\footnotesize(300k lignes)};

  \node[dim, above=of fr] (dd) {dim\_date\\\footnotesize(calendrier)};
  \node[dim, below=of fr] (du) {dim\_users\\\footnotesize(10k abonnés)};
  \node[dim, left=of fr]  (dz) {dim\_zones\\\footnotesize(23 zones)};
  \node[dim, above=of fc] (dt) {dim\_tranches\\\footnotesize(grille 2026)};

  \draw[arr] (fr) -- (dd);
  \draw[arr] (fr) -- (du);
  \draw[arr] (fr) -- (dz);
  \draw[arr] (fc) -- (dd);
  \draw[arr] (fc) -- (du);
  \draw[arr] (fc) -- (dt);
\end{tikzpicture}
\caption{Schéma en étoile de l'entrepôt de données Woyofal}
\label{fig:dwh}
\end{figure}

\subsection{Tables clés}

\paragraph{dim\_users} 10~000 abonnés avec type compteur (DPP/PPP), zone géographique, consommation quotidienne moyenne et objectif mensuel.

\paragraph{dim\_tranches} Table de référence des tarifs 2026 insérée statiquement à la création du schéma. Contient les seuils, prix par kWh et description pour DPP et PPP.

\paragraph{fact\_recharges} Enregistre chaque transaction (montant, kWh obtenus, tranche atteinte, redevance). Clés étrangères vers \texttt{dim\_date}, \texttt{dim\_users}, \texttt{dim\_zones}.

\paragraph{fact\_consumption} Enregistre la consommation quotidienne (kWh, coût FCFA, tranche, cumul mensuel, économie baisse 10\%). Clés étrangères vers \texttt{dim\_date}, \texttt{dim\_users}, \texttt{dim\_tranches}.

\subsection{Optimisations}

\begin{itemize}
  \item Index sur toutes les clés étrangères et colonnes fréquemment filtrées (\texttt{mois}, \texttt{annee}, \texttt{type\_compteur}) ;
  \item Vues matérialisées pour les agrégats BI fréquents ;
  \item Data marts dédiés (\texttt{sql/05\_create\_data\_marts.sql}) pour les requêtes Streamlit/PowerBI ;
  \item Commande \texttt{COPY} server-side pour l'ingestion des gros CSV (évite les problèmes d'encodage).
\end{itemize}

\section{Transformations dbt}

Les transformations métier sont versionnées et testées avec \textbf{dbt} (\textit{data build tool}). Le projet \texttt{woyofal\_dbt/} contient :

\begin{itemize}
  \item \textbf{Modèles staging} — nettoyage et standardisation des tables brutes ;
  \item \textbf{Modèles intermédiaires} — calcul des agrégats mensuels par utilisateur et par zone ;
  \item \textbf{Modèles marts} — tables prêtes pour les dashboards ;
  \item \textbf{Tests dbt} — vérification unicité des clés, non-nullité, cohérence référentielle.
\end{itemize}

%% ============================================================
\chapter{Analyse Exploratoire et Machine Learning}
\label{ch:eda_ml}
%% ============================================================

\section{Analyse Exploratoire des Données (EDA)}

L'EDA a été conduite sur l'ensemble des 330~000 enregistrements via le script \texttt{scripts/run\_eda.py} et les notebooks Jupyter \texttt{01\_EDA\_Initiale.ipynb} et \texttt{02\_EDA\_Complete.ipynb}.

\subsection{Distribution de la consommation}

\begin{infobox}[Principaux indicateurs EDA]
\begin{itemize}
  \item Consommation quotidienne moyenne : \textbf{$\approx$ 5,2~kWh/jour}
  \item Montant de recharge moyen : \textbf{$\approx$ 8~400~FCFA}
  \item Part des recharges restant en T1 : \textbf{68\%}
  \item 73\% des journées de consommation se situent en tranche~1
  \item Les zones urbaines consomment en moyenne \textbf{1,8$\times$} plus que les zones rurales
  \item Saisonnalité marquée : hausse en saison sèche (décembre--mars)
\end{itemize}
\end{infobox}

\subsection{Corrélation montant--kWh}

L'analyse de corrélation entre le montant rechargé (FCFA) et les kWh obtenus révèle une \textbf{corrélation non-linéaire} due à la progressivité des tranches. À montant égal, un utilisateur avec un cumul élevé obtient moins de kWh qu'un utilisateur en début de mois — valeur pédagogique directe pour la plateforme.

\subsection{Évolution mensuelle}

L'analyse de la tendance mensuelle (2024--2026) montre un point d'inflexion post-janvier 2026 correspondant à l'entrée en vigueur de la nouvelle grille tarifaire (baisse 10\% T1).

\section{Modèle de Machine Learning}

\subsection{Objectif du modèle}

Le modèle répond à la question : \textbf{étant donné le profil d'un utilisateur et son historique du mois en cours, va-t-il dépasser la tranche~1 lors de sa prochaine recharge ?}

Cette prédiction permet d'alerter proactivement l'utilisateur avant qu'il ne passe en tranche plus coûteuse.

\subsection{Préparation des données}

Le dataset ML (\texttt{data/processed/dataset\_ml\_ready.csv}) est construit à partir de la table \texttt{fact\_consumption} enrichie par des features ingénieurs :

\begin{itemize}
  \item Cumul mensuel courant, moyenne des 7 derniers jours, écart-type de consommation ;
  \item Jour du mois, saison, type de zone (urbain/rural), type compteur (DPP/PPP) ;
  \item Montant moyen des recharges passées, fréquence de recharge.
\end{itemize}

La variable cible \texttt{target\_depasse\_t1} est un booléen indiquant si le cumul mensuel dépasse 150~kWh. Le split train/test est stratifié (80\%/20\%).

\subsection{Algorithme : Random Forest}

Un \textbf{Random Forest Classifier} (scikit-learn) a été sélectionné pour sa robustesse aux données tabulaires et sa capacité à capturer les interactions non-linéaires.

\begin{codebox}[Configuration du modèle]
\begin{lstlisting}[language=Python, numbers=none]
clf = RandomForestClassifier(
    n_estimators=100,
    max_depth=15,
    min_samples_split=50,
    random_state=42,
    n_jobs=-1          # Parallélisation sur tous les cœurs
)
\end{lstlisting}
\end{codebox}

\subsection{Résultats}

\begin{table}[h!]
\centering
\caption{Métriques du modèle Random Forest — Prédiction dépassement T1}
\label{tab:ml}
\begin{tabularx}{0.65\textwidth}{l r}
\toprule
\rowcolor{marine!10}
Métrique & Valeur \\
\midrule
Accuracy  & \textbf{99,75\%} \\
Precision & 99,80\% \\
Recall    & 99,77\% \\
F1-Score  & 99,79\% \\
ROC-AUC   & \textbf{99,999\%} \\
\midrule
Jeu de test & 20\% de 330~000 lignes \\
Durée d'entraînement & $<$ 30 secondes \\
\bottomrule
\end{tabularx}
\end{table}

\begin{infobox}[Interprétation des résultats]
Les métriques très élevées s'expliquent par la nature déterministe des données (les seuils de tranche sont des règles exactes intégrées dans la génération). Le modèle apprend la règle tarifaire. Sa valeur réside dans sa capacité à \textbf{anticiper} le dépassement \textit{avant} la recharge, à partir de l'historique partiel du mois, permettant des alertes préventives.
\end{infobox}

\subsection{Feature Importance}

Les 3 features les plus déterminantes :
\begin{enumerate}
  \item \texttt{conso\_cumul\_mois} — cumul mensuel courant (principal prédicteur) ;
  \item \texttt{conso\_kwh\_mean\_7j} — consommation moyenne des 7 derniers jours ;
  \item \texttt{jour\_du\_mois} — le risque de dépassement augmente en fin de mois.
\end{enumerate}

%% ============================================================
\chapter{Implémentation}
\label{ch:impl}
%% ============================================================

\section{Moteur de calcul tarifaire}

Le principal défi technique est l'implémentation d'un moteur de calcul conforme à la grille tarifaire progressive de Senelec. Ce moteur est exécuté entièrement côté client (navigateur), sans appel à un serveur, garantissant la disponibilité hors connexion.

\subsection{Algorithme de simulation directe (FCFA $\rightarrow$ kWh)}

\begin{codebox}[Algorithme — Mode direct]
\begin{lstlisting}[language=JavaScript, numbers=none]
const simulerRechargeLocal = (montant, typeCompteur, cumulActuel, avecRedevance) => {
  const redevance  = avecRedevance ? 429 : 0
  const taxe       = Math.round(montant * 0.025)   // 2.5% taxe communale
  let montantReste = montant - redevance - taxe

  let cumul    = cumulActuel
  let kwhTotal = 0

  for (let t = 1; t <= 4; t++) {
    if (montantReste <= 0) break
    const disponible = maxTranche ? maxTranche - cumul : Infinity
    const kwh        = Math.min(montantReste / prix[t], disponible)
    kwhTotal    += kwh
    cumul       += kwh
    montantReste -= kwh * prix[t]
  }

  return { kwh_obtenus: kwhTotal, tranche_finale: ..., detail_tranches: ... }
}
\end{lstlisting}
\end{codebox}

\textbf{Exemple concret} : pour 10~000~FCFA avec cumul~=~0 et sans redevance :
\begin{align*}
\text{Taxe}          &= 10\,000 \times 2{,}5\% = 250~\text{FCFA} \\
\text{Net}           &= 10\,000 - 250 = 9\,750~\text{FCFA} \\
\text{T1 (82 F/kWh)} &= \min\!\left(\frac{9750}{82},\;150\right) = 118{,}9~\text{kWh} \\
\text{Total}         &= 118{,}9~\text{kWh en Tranche 1}
\end{align*}

\subsection{Algorithme de calcul inverse (kWh $\rightarrow$ FCFA)}

Le calcul inverse est une fonctionnalité exclusive : l'utilisateur saisit le nombre de kWh souhaités et la plateforme calcule le montant exact à recharger.

\begin{codebox}[Algorithme — Mode inverse]
\begin{lstlisting}[language=JavaScript, numbers=none]
const calculerMontantPourKwh = (kwhVoulus, cumulActuel, typeCompteur, avecRedevance) => {
  const tarifs    = TARIFS_LOCAL[typeCompteur]
  const redevance = avecRedevance ? 429 : 0
  let cumul       = cumulActuel
  let montantNet  = 0
  let kwhReste    = kwhVoulus

  for (let t = 1; t <= 3; t++) {
    if (kwhReste <= 0) break
    const disponible      = Math.max(0, maxTranche - Math.max(cumul, minTranche))
    const kwhDansTranche  = Math.min(kwhReste, disponible)
    montantNet += kwhDansTranche * tarifs[t].prix
    cumul      += kwhDansTranche
    kwhReste   -= kwhDansTranche
  }

  const montantAvantTaxe = montantNet + redevance
  const montantBrut      = montantAvantTaxe / (1 - 0.025)   // Inversion taxe
  return Math.ceil(montantBrut / 100) * 100                 // Arrondi au 100F sup.
}
\end{lstlisting}
\end{codebox}

\section{Interface utilisateur}

\subsection{Structure du frontend}

\begin{lstlisting}[language={},caption={Arborescence du frontend React}]
frontend-react/
├── src/
│   ├── components/
│   │   ├── HomePage.jsx               # Page d'accueil + mini-calc
│   │   ├── SimulateurRecharge.jsx     # Simulateur complet
│   │   ├── Dashboard.jsx              # Espace utilisateur (KPIs)
│   │   ├── AIChat.jsx                 # Assistant IA flottant
│   │   ├── AuthModal.jsx              # Authentification JWT
│   │   ├── GuideTarifs.jsx            # Grille tarifaire DPP/PPP
│   │   ├── FAQ.jsx                    # Questions fréquentes
│   │   ├── Conseils.jsx               # Calculateur d'économies
│   │   └── HistoriqueConsommation.jsx # Historique [Protégé]
│   ├── context/
│   │   └── AuthContext.jsx            # État global auth
│   ├── services/
│   │   └── api.js                     # Client Axios (dev only)
│   ├── App.jsx                        # Routeur principal
│   └── index.css
├── tailwind.config.js                 # Thème couleurs
└── vercel.json                        # Rewrite SPA
\end{lstlisting}

\subsection{Page d'accueil}

La page d'accueil (\texttt{HomePage.jsx}) est organisée en 6 sections : Hero avec mini-calculateur interactif en temps réel, À propos, Tarifs DPP 2026, Comment ça marche, FAQ accordéon, et CTA final. Le mini-calculateur visualise les kWh obtenus pour un glisseur de montant, avec affichage de la tranche atteinte (T1~vert, T2~orange, T3~rouge).

\subsection{Simulateur complet}

Le simulateur (\texttt{SimulateurRecharge.jsx}) offre deux modes :
\begin{itemize}
  \item \textbf{Mode FCFA $\rightarrow$ kWh} (direct) : saisie du montant, du cumul, du type de compteur et de la redevance ; résultat détaillé par tranche.
  \item \textbf{Mode kWh $\rightarrow$ FCFA} (inverse, exclusif) : saisie des kWh souhaités ; retourne le montant exact à recharger.
\end{itemize}

\subsection{Assistant IA}

Le composant \texttt{AIChat.jsx} intègre un chatbot flottant basé sur l'API \textbf{OpenRouter} avec une cascade de six modèles : Llama~3.3~70B $\rightarrow$ Gemma~3 $\rightarrow$ Qwen~3 $\rightarrow$ Mistral~7B, avec mécanisme de fallback local intelligent. Il comprend les questions en langage naturel, effectue les calculs tarifaires à la volée et gère l'historique de conversation.

\section{Déploiement CI/CD}

\begin{lstlisting}[language={},caption={vercel.json — Configuration de déploiement SPA}]
{
  "buildCommand": "npm run build",
  "outputDirectory": "dist",
  "rewrites": [
    { "source": "/(.*)", "destination": "/index.html" }
  ]
}
\end{lstlisting}

La règle \textit{rewrite} est essentielle pour React Router : toutes les URLs sont redirigées vers \texttt{index.html} pour être gérées côté client. Le workflow GitHub Actions déclenche automatiquement le déploiement Vercel à chaque \texttt{git push origin main}.

%% ============================================================
\chapter{Résultats et Tests}
\label{ch:resultats}
%% ============================================================

\section{Fonctionnalités livrées}

\begin{table}[h!]
\centering
\caption{Bilan des fonctionnalités}
\label{tab:fonctionnalites}
\begin{tabularx}{\textwidth}{l X c}
\toprule
\rowcolor{marine!10}
Fonctionnalité & Description & Statut \\
\midrule
Pipeline ETL Airflow   & DAG quotidien, ingest 330k lignes PostgreSQL & \textcolor{green!60!black}{\textbf{✓}} \\
Data Warehouse         & Star schema, 4 dims, 2 facts, data marts     & \textcolor{green!60!black}{\textbf{✓}} \\
Transformations dbt    & Modèles staging, intermédiaires, marts        & \textcolor{green!60!black}{\textbf{✓}} \\
EDA complète           & Distribution, corrélations, saisonnalité     & \textcolor{green!60!black}{\textbf{✓}} \\
Modèle ML              & RandomForest 99,75\% accuracy                & \textcolor{green!60!black}{\textbf{✓}} \\
Mini-calculateur       & Glisseur FCFA, 4 tranches, tranche atteinte & \textcolor{green!60!black}{\textbf{✓}} \\
Simulateur direct      & FCFA $\rightarrow$ kWh détaillé par tranche & \textcolor{green!60!black}{\textbf{✓}} \\
Simulateur inverse     & kWh $\rightarrow$ FCFA (exclusif)            & \textcolor{green!60!black}{\textbf{✓}} \\
Support DPP/PPP        & Les deux types de compteur                  & \textcolor{green!60!black}{\textbf{✓}} \\
Historique localStorage & Sauvegarde sans compte                     & \textcolor{green!60!black}{\textbf{✓}} \\
Auth JWT               & Connexion / Inscription                     & \textcolor{green!60!black}{\textbf{✓}} \\
Assistant IA           & OpenRouter multi-modèles, fallback local    & \textcolor{green!60!black}{\textbf{✓}} \\
Responsive mobile      & Testé iOS + Android                         & \textcolor{green!60!black}{\textbf{✓}} \\
Déploiement Vercel     & Production en ligne                         & \textcolor{green!60!black}{\textbf{✓}} \\
CI/CD GitHub           & Push $\rightarrow$ déploiement auto         & \textcolor{green!60!black}{\textbf{✓}} \\
\bottomrule
\end{tabularx}
\end{table}

\section{Validation du moteur de calcul}

\begin{table}[h!]
\centering
\caption{Tests de validation du simulateur DPP}
\label{tab:tests}
\begin{tabularx}{\textwidth}{r r r r X}
\toprule
\rowcolor{marine!10}
Montant & Cumul & Redevance & Résultat (kWh) & Tranche finale \\
\midrule
5~000 F  & 0 & Non & 60,9 kWh              & T1 \\
10~000 F & 0 & Non & 118,9 kWh             & T1 \\
15~000 F & 0 & Non & 150 (T1) + 10,1 (T2) & T2 \\
20~000 F & 0 & Non & 150 (T1) + 62,8 (T2) & T2 \\
5~000 F  & 0 & Oui & 55,7 kWh              & T1 \\
\midrule
\textit{Inverse :} 150 kWh & 0  & Non & 12~300 F & T1 (seuil exact) \\
\textit{Inverse :} 100 kWh & 50 & Non & 8~200 F  & T2 \\
\bottomrule
\end{tabularx}
\end{table}

\section{Bugs identifiés et résolus}

\begin{infobox}[Corrections apportées en cours de développement]
\begin{enumerate}
  \item \textbf{Page blanche (prod)} — Variable \texttt{tarifs} référencée dans le JSX après suppression de son état React ; correction : suppression du bloc JSX orphelin.
  \item \textbf{Appel \texttt{localhost:5000} sur mobile} — Le simulateur appelait l'API locale en production ; correction : moteur de calcul déplacé entièrement côté client.
  \item \textbf{\texttt{saveSimulation} manquant} — Fonction supprimée des imports mais encore appelée ; correction : remplacement par sauvegarde \texttt{localStorage}.
  \item \textbf{Secret Vercel manquant} — \texttt{vercel.json} référençait \texttt{@api\_url\_prod} inexistant ; correction : suppression de la clé \texttt{env}.
\end{enumerate}
\end{infobox}

%% ============================================================
\chapter{Défis et Leçons Apprises}
%% ============================================================

\section{Défis rencontrés}

\subsection{Calcul tarifaire progressif}
La grille Senelec ne s'applique pas sur le montant total mais \textit{progressivement} : chaque FCFA est d'abord converti aux kWh de la tranche courante, jusqu'au seuil, puis la tranche suivante prend le relais. Cet algorithme itératif a nécessité une attention particulière pour les cas limites (cumul déjà en T2, montant insuffisant pour remplir T1, inversion de la taxe en mode inverse).

\subsection{Indépendance du backend en production}
Le simulateur appelait initialement une API Node.js locale (\texttt{localhost:5000}). En production sur Vercel, aucun backend n'étant disponible, cela provoquait une page blanche sur mobile. La solution a été de dupliquer le moteur de calcul en JavaScript pur, entièrement dans le navigateur.

\subsection{Volume de données et performance ETL}
L'ingestion de 300~000 lignes via des insertions unitaires atteignait des délais prohibitifs. L'adoption du \texttt{COPY} server-side PostgreSQL et du batch processing (tailles 5~000 / 2~000) a ramené le temps ETL total à 8--12 minutes.

\subsection{Cohérence visuelle multi-composants}
Avec 10 composants React partageant la même identité visuelle, la centralisation des couleurs dans \texttt{tailwind.config.js} (\texttt{primary: '\#DC2626'}, \texttt{accent: '\#1A1D2E'}) a maintenu la cohérence.

\section{Leçons apprises}

\begin{itemize}
  \item \textbf{Architecture offline-first} : concevoir le calcul métier côté client dès le départ évite les dépendances runtime en production ;
  \item \textbf{Test sur mobile réel} : le rendu sur desktop ne suffit pas — les tests sur téléphone ont révélé les bugs les plus critiques ;
  \item \textbf{Séparation des préoccupations} : un moteur de calcul pur (sans React) est plus facile à tester et réutilisable (CLI, API, frontend) ;
  \item \textbf{COPY server-side} : pour les gros volumes CSV, l'ingestion via le serveur PostgreSQL est incomparablement plus rapide que les insertions Python row-by-row ;
  \item \textbf{dbt pour la traçabilité} : versionner les transformations SQL permet d'auditer le lineage et d'éviter les transformations silencieuses.
\end{itemize}

%% ============================================================
\chapter{Conclusion et Perspectives}
\label{ch:conclusion}
%% ============================================================

\section{Bilan}

Ce projet a permis de concevoir et déployer une \textbf{plateforme de données complète} couvrant l'intégralité de la chaîne, depuis la génération des données jusqu'à l'interface utilisateur en production. Tous les objectifs fixés ont été atteints :

\begin{itemize}
  \item[$\checkmark$] Pipeline ETL Airflow opérationnel, ingestion de 330~000 lignes dans PostgreSQL ;
  \item[$\checkmark$] Entrepôt de données star schema avec 4 dimensions, 2 tables de faits, data marts et vues matérialisées optimisées ;
  \item[$\checkmark$] Transformations dbt versionnées et testées ;
  \item[$\checkmark$] EDA complète (distribution, corrélations, saisonnalité, analyse tarifaire) ;
  \item[$\checkmark$] Modèle Random Forest avec 99,75\% d'accuracy (prédiction dépassement T1) ;
  \item[$\checkmark$] Moteur de calcul DPP/PPP 2026 implémenté et validé, avec mode inverse exclusif ;
  \item[$\checkmark$] Interface React responsive déployée en production sur Vercel avec CI/CD.
\end{itemize}

La plateforme est accessible publiquement à l'adresse : \\
\url{https://frontend-react-wine.vercel.app}

\section{Limites actuelles}

\begin{itemize}
  \item L'assistant IA nécessite une connexion internet (API OpenRouter) ;
  \item L'historique des recharges est stocké en \texttt{localStorage} uniquement — non synchronisé entre appareils ;
  \item Le backend Node.js mock n'est pas déployé : authentification et dashboard ne fonctionnent qu'en développement local.
\end{itemize}

\section{Perspectives}

\textbf{Court terme (1--3 mois) :}
\begin{itemize}
  \item Déploiement du backend (Railway/Render) pour activer l'authentification en production ;
  \item Exposition du modèle ML via une API REST pour déclencher des alertes de dépassement en temps réel ;
  \item Export PDF des simulations depuis l'interface.
\end{itemize}

\textbf{Moyen terme (3--6 mois) :}
\begin{itemize}
  \item Connexion à des données réelles Senelec (si ouverture des API) pour remplacer les données synthétiques ;
  \item Intégration d'un modèle de prévision de consommation (Prophet/LightGBM) pour anticiper la facture mensuelle ;
  \item Dashboard PowerBI/Streamlit accessible aux gestionnaires de zones ;
  \item Application mobile React Native avec notifications push d'alerte de tranche.
\end{itemize}

%% ============================================================
%% ANNEXES
%% ============================================================
\appendix

\chapter{Captures d'écran}
\label{ann:screenshots}

\textit{Insérer les captures ici, exemple :}
\begin{lstlisting}[language={},numbers=none]
\begin{figure}[h!]
  \centering
  \includegraphics[width=0.9\textwidth]{screenshots/homepage.png}
  \caption{Page d'accueil — section Hero avec mini-calculateur}
\end{figure}
\end{lstlisting}

\chapter{Extraits de code clés}
\label{ann:code}

\section*{Calcul inverse complet}

\begin{lstlisting}[language=JavaScript, caption={calculerMontantPourKwh — calcul inverse DPP/PPP}]
const calculerMontantPourKwh = (kwhVoulus, cumulActuel, typeCompteur, avecRedevance) => {
  const tarifs    = TARIFS_LOCAL[typeCompteur]
  const redevance = avecRedevance ? 429 : 0
  let cumul       = cumulActuel
  let montantNet  = 0
  let kwhReste    = kwhVoulus

  for (let t = 1; t <= 3; t++) {
    if (kwhReste <= 0) break
    const maxTranche  = tarifs[t].max
    const minTranche  = t === 1 ? 0 : tarifs[t - 1].max
    const disponible  = maxTranche
      ? Math.max(0, maxTranche - Math.max(cumul, minTranche))
      : Infinity
    if (disponible <= 0) continue
    const kwhDansTranche = Math.min(kwhReste, disponible)
    montantNet += kwhDansTranche * tarifs[t].prix
    cumul      += kwhDansTranche
    kwhReste   -= kwhDansTranche
  }

  const montantAvantTaxe = montantNet + redevance
  const montantBrut      = montantAvantTaxe / (1 - 0.025)
  return Math.ceil(montantBrut / 100) * 100
}
\end{lstlisting}

\chapter{Liens et ressources}
\label{ann:liens}

\begin{tabular}{ll}
\toprule
Ressource & URL \\
\midrule
Site en production & \url{https://frontend-react-wine.vercel.app} \\
Dépôt GitHub & \url{https://github.com/ndoye09/woyofal2026} \\
Décision CRSE 2025-140 & CRSE Sénégal — grille tarifaire Senelec 2026 \\
\bottomrule
\end{tabular}

\end{document}
